\chapter{Limits and troubleshooting}\label{ch:limits}

\section{Limits}

\begin{description}
  \item[Experimental.] miniVivisection relies on internal code of Rhino~8's Script Editor (RhinoCode 8.36 when it was
        written). A Rhino service release can change it. The plug-in then follows nothing; it doesn't fail.
  \item[Python~3 first.] Inline values, tooltips with scope, the Python prompt in the paused frame and Edit in Rhino are
        for Python~3 scripts only.
  \item[IronPython~2] (\code{\#! python2}): FreePause, step over across modules, \vn{Watch} drawing and baking of plain variable
        names, the hover preview and tooltip values of plain names, the value tree, conditional breakpoints, hit counts and
        logpoints work. Left out, quietly: expressions in watches and the Watch tab, scoped tooltips, values at the ends of
        lines, “was”, Edit in Rhino and Auto edit, visualizers and Interpret as for raw Python data, the Python prompt and
        Claude's evaluation in the paused frame. While one is debugged, an amber notice on \vn{Python} and \vn{Claude} says so.
        Full IronPython~2 support is planned.
  \item[C\#:] Watches of plain variable names and conditional breakpoints with a small evaluator; largely untested.
  \item[Two watch lists.] The editor's Watch tab shows values; \vn{Watch} draws. Each has its own list.
  \item[Grasshopper:] Python~3 script components through their own script editor. The legacy GhPython (IronPython~2)
        component isn't covered. With a component's editor closed and no hidden editor, \keys{Ctrl+Alt+B} can only stop it,
        not pause it.
  \item[Not traced:] Python code run by .NET threads.
\end{description}

\section{Troubleshooting}

\begin{xltabular}{\linewidth}{@{}P{0.34\linewidth}L@{}}
\toprule
\textsf{\textbf{What happens}} & \textsf{\textbf{What to check}} \\ \midrule
Nothing is drawn, no values in the code & Is the script paused? Is \vn{Debug} ON in the top bar? Is \setting{Draw watches in Rhino} on? Is the watch's eye open? \\
\key{F5} doesn't stop                   & Without a breakpoint, \key{F5} is a plain run. Use \ui{Debug} (\see{sec:runbuttons}). \\
\keys{Ctrl+Alt+B} does nothing          & An error sound means nothing could be paused or stopped; \vn{Log} says why. A plain run with \setting{RunPythonScript command} on \emph{Normal}, or a Grasshopper component in \emph{Normal} mode, is left alone. \\
A breakpoint in a thread doesn't pause  & The freeze guard skipped it: the script runs on Rhino's UI thread. Add \code{\# async:true} (\see{sec:freezeguard}). \\
A condition always pauses               & It fails: a typo or a missing name. Its marker is red, and \vn{Log} says why. \\
The editor of a paused script is gone   & \code{miniVivisection StopPaused}. \\
Claude says “the rhino MCP didn't answer” & RhinoAI's server isn't running in this Rhino: run \code{MCPStart}, or turn on \setting{Start RhinoAI's server for Claude} (\see{sec:rhinoai}). \\
Claude doesn't work in the model        & RhinoAI isn't installed, or its server isn't running; or the tool is a destructive one: \code{/allow on}. \\
Grasshopper shows “expired during a solution” & Something recomputed while a component was paused; it should be held back (\see{sec:ghheld}). Report it with the detailed log. \\
A Grasshopper component can't be deleted & A component is paused; deleting waits for the pause to end. \\
After a Rhino update, a feature is gone & The Script Editor's internal code may have changed. Look at the detailed log, and for a newer miniVivisection. \\
\bottomrule
\end{xltabular}

\section{Reporting a problem}

Note the version and build (the top bar, or the load message), what you did, and what happened. Open the log folder
(\code{miniVivisection LogFolder}) and attach \file{minivivisection.log}, or \file{minivivisection.crashed-(time).log} after a crash. For
a problem that is hard to see, set \setting{Detailed log} to \emph{Debug} or \emph{Trace} for the area concerned first.
